% =============================================================================
%  Notas de clase — Sesión 23: Argumentos de consola
%  Programación Avanzada — Universidad Panamericana
%  Compilar: pdflatex sesion23_argumentos_consola.tex
% =============================================================================
\documentclass[12pt, oneside]{article}
\usepackage{progavanzada}
\usepackage{array}

\begin{document}

\portada{Sesión 23}{Argumentos de consola}{2026}

\tableofcontents
\newpage

% =============================================================================
\section{1969: ``haz una sola cosa y hazla bien''}
% =============================================================================

En 1969, Ken Thompson diseñó la primera versión de Unix en un PDP-7
que nadie más usaba en los Laboratorios Bell.  Dennis Ritchie se unió
poco después, y juntos establecieron una filosofía de diseño que sigue
siendo vigente más de cincuenta años después:

\begin{center}
\textit{``Escribe programas que hagan una sola cosa y la hagan bien.''}\\
\textit{``Escribe programas que puedan trabajar juntos.''}
\end{center}

Para que los programas trabajen juntos necesitan dos mecanismos: una
forma de recibir instrucciones desde el exterior, y una forma de pasar
sus resultados a otro programa.  La segunda se resolvió con las
\textbf{tuberías} (\textit{pipes}); la primera, con los
\textbf{argumentos de consola}.

Cuando escribes:

\begin{verbatim}
$ g++ -O2 -o hola hola.cpp
\end{verbatim}

el shell toma esa línea, la separa en cinco cadenas y se las pasa al
compilador.  El compilador las recibe como un arreglo de cadenas llamado
\cpp{argv} (de \textit{argument vector}) y un entero llamado \cpp{argc}
(de \textit{argument count}).  Llevas un semestre usando \texttt{g++}
sin saber que le estás pasando cuatro argumentos por \cpp{argv}.

Esta convención fue diseñada por Dennis Ritchie al desarrollar el
lenguaje~C alrededor de 1972 y quedó formalizada en el estándar
ANSI~C de 1989.  Hoy es parte del estándar de C++ y de prácticamente
todo lenguaje compilado moderno.

\begin{tecnico}[La filosofía Unix y las tuberías]
  Los programas \texttt{grep}, \texttt{wc}, \texttt{sort} y \texttt{ls}
  son extremadamente pequeños porque cada uno hace una sola cosa.  Su
  poder real proviene de combinarlos:
  \begin{verbatim}
$ cat notas.txt | grep "error" | wc -l
  \end{verbatim}
  Este comando cuenta cuántas líneas de \texttt{notas.txt} contienen
  la palabra ``error''.  Cada programa recibe su configuración por
  \cpp{argv} y sus datos por la entrada estándar.
  \par\smallskip
  \textit{Referencia:} Ritchie, D.\,M., \& Thompson, K.~(1974).
  ``The UNIX time-sharing system.''
  \textit{Communications of the ACM}, 17(7), 365--375.
\end{tecnico}

% =============================================================================
\section{\texttt{main} con parámetros: \texttt{argc} y \texttt{argv}}
% =============================================================================

Hasta ahora siempre has escrito:

\begin{codigo}
int main(){
    ...
}
\end{codigo}

Esa forma no recibe argumentos.  La versión completa de \cpp{main} es:

\begin{codigo}
int main(int argc, char* argv[]){
    ...
}
\end{codigo}

\begin{description}
  \item[\cpp{argc}]
    Número de cadenas recibidas, incluyendo el nombre del programa.
    Su valor mínimo es~1.
  \item[\cpp{argv}]
    Arreglo de cadenas estilo~C.  Cada elemento es un \cpp{char*}
    que apunta a una cadena terminada en \texttt{'\textbackslash{}0'}.
    \begin{itemize}
      \item \cpp{argv[0]}: nombre o ruta del ejecutable.
      \item \cpp{argv[1]} a \cpp{argv[argc-1]}: los argumentos del usuario.
    \end{itemize}
\end{description}

Ejemplo: si ejecutas:

\begin{verbatim}
$ ./mi_programa hola mundo 42
\end{verbatim}

entonces:

\begin{center}
\begin{tabular}{cll}
  \toprule
  \textbf{Variable}  & \textbf{Valor}            & \textbf{Descripción} \\
  \midrule
  \cpp{argc}         & \texttt{4}                & cuatro cadenas en total \\
  \cpp{argv[0]}      & \texttt{"./mi\_programa"} & nombre del ejecutable \\
  \cpp{argv[1]}      & \texttt{"hola"}           & primer argumento \\
  \cpp{argv[2]}      & \texttt{"mundo"}          & segundo argumento \\
  \cpp{argv[3]}      & \texttt{"42"}             & tercer argumento \\
  \bottomrule
\end{tabular}
\end{center}

\begin{advertencia}
  \cpp{argv[3]} es la cadena \texttt{"42"}, no el número \texttt{42}.
  Si quieres usarlo en una operación aritmética hay que convertirlo
  explícitamente.  Este es uno de los errores más comunes al trabajar
  con argumentos de consola.
\end{advertencia}

\begin{recuerda}[Dos notaciones equivalentes para \texttt{argv}]
  Encontrarás las siguientes dos formas en el código; ambas son
  idénticas para el compilador:
  \begin{itemize}
    \item \cpp{char* argv[]} — arreglo de apuntadores a \cpp{char}
    \item \cpp{char** argv} — apuntador a apuntador a \cpp{char}
  \end{itemize}
  Elige la que te resulte más clara.  En este curso usaremos
  \cpp{char** argv} en los ejemplos cortos y \cpp{char* argv[]}
  cuando queremos enfatizar que es un arreglo.
\end{recuerda}

% -----------------------------------------------------------------------------
\subsection{Primer programa: contar argumentos}
% -----------------------------------------------------------------------------

El programa más sencillo que usa \cpp{argc}/\cpp{argv} simplemente
reporta cuántos argumentos recibió:

\begin{codigo}[caption=argumentos1.cpp]
#include <iostream>
using namespace std;

int main(int argc, char **argv){
    cout << "Recibí " << argc << " argumentos\n";
}
\end{codigo}

Ejecución en terminal:
\begin{verbatim}
$ ./argumentos1 hola mundo
Recibí 3 argumentos
\end{verbatim}

Observa que el ejecutable mismo cuenta como un argumento, por eso el
resultado es~3 y no~2.

% -----------------------------------------------------------------------------
\subsection{Iterar e imprimir los argumentos}
% -----------------------------------------------------------------------------

Para mostrar los argumentos que el usuario proporcionó (sin el nombre
del programa) iteramos desde \cpp{i=1}:

\begin{codigo}[caption=argumentos2.cpp]
#include <iostream>
using namespace std;

int main(int argc, char **argv){
    cout << "Recibí " << argc << " argumentos\n";
    for (int i = 1; i < argc; i++){
        cout << "\t" << argv[i] << endl;
    }
}
\end{codigo}

Ejecución en terminal:
\begin{verbatim}
$ ./argumentos2 rojo verde azul
Recibí 4 argumentos
        rojo
        verde
        azul
\end{verbatim}

\begin{recuerda}
  El shell divide la línea de comandos usando los espacios como
  separadores.  Si quieres pasar una cadena con espacios como un solo
  argumento, enciérrala en comillas:
  \begin{verbatim}
$ ./argumentos2 "hola mundo"
  \end{verbatim}
  En ese caso \cpp{argc} es~2 y \cpp{argv[1]} es \texttt{"hola mundo"}
  (una sola cadena, no dos).
\end{recuerda}

% -----------------------------------------------------------------------------
\subsection{\texttt{argv[i]} es una cadena de C: recorrerla carácter por carácter}
% -----------------------------------------------------------------------------

Cada elemento de \cpp{argv} es un \cpp{char*}: un apuntador al primer
carácter de una cadena terminada en \texttt{'\textbackslash{}0'}.  Eso
significa que puedes recorrerla carácter a carácter exactamente igual
que cualquier cadena de C:

\begin{codigo}[caption=argumentos4.cpp]
#include <iostream>
using namespace std;

int main(int argc, char **argv){
    cout << "Recibí " << argc << " argumentos\n";
    char *x = argv[1];
    int i = 0;
    while (true){
        if (x[i] == '\0') break;
        cout << x[i] << endl;
        i++;
    }
}
\end{codigo}

Ejecución en terminal:
\begin{verbatim}
$ ./argumentos4 hola
Recibí 2 argumentos
h
o
l
a
\end{verbatim}

¿Por qué termina en \texttt{'\textbackslash{}0'}?  Porque el shell
construye cada argumento como una cadena de C estándar.  El bucle
\cpp{while(true)} con ruptura al encontrar el centinela es el patrón
clásico para recorrer cadenas de C.

\begin{consejo}
  Este ejemplo conecta directamente con lo que ya sabes sobre
  apuntadores y arreglos.  \cpp{argv} es un \cpp{char**}: un apuntador
  que apunta a otro apuntador.  \cpp{argv[1]} devuelve el segundo
  \cpp{char*}, y \cpp{argv[1][0]} el primer carácter de ese argumento.
\end{consejo}

% =============================================================================
\section{Verificar el número de argumentos}
% =============================================================================

Los programas robustos siempre verifican que recibieron el número
correcto de argumentos antes de intentar usarlos.  Acceder a
\cpp{argv[2]} cuando \cpp{argc} es~2 es comportamiento indefinido;
en la práctica provoca un \textit{segmentation fault} o lee memoria
basura.

El siguiente programa espera exactamente dos números; si no los recibe,
muestra un mensaje de uso y termina:

\begin{codigo}[caption=argumentos3.cpp]
#include <iostream>
using namespace std;

int main(int argc, char **argv){
    if (argc != 3){
        cerr << "Debes mandar dos números\n";
        cerr << "./suma n1 n2\n";
        return 1;
    }
    int x, y;
    x = atoi(argv[1]);
    y = atoi(argv[2]);
    cout << "La suma de " << x << " + " << y << " = " << x + y << endl;
}
\end{codigo}

Ejecución en terminal:
\begin{verbatim}
$ ./argumentos3 10 25
La suma de 10 + 25 = 35

$ ./argumentos3 10
Debes mandar dos números
./suma n1 n2
\end{verbatim}

Dos detalles importantes:

\begin{enumerate}
  \item El mensaje de error va a \cpp{cerr} (la salida de error estándar),
        no a \cpp{cout}.  Así el usuario puede redirigir la salida normal
        sin mezclarla con mensajes de error.

  \item El programa devuelve \cpp{1} para indicar al shell que algo salió
        mal.  Por convención, \cpp{return 0} significa éxito y cualquier
        otro valor indica error.
\end{enumerate}

\begin{consejo}
  Usa \cpp{argv[0]} en el mensaje de uso en lugar de escribir el nombre
  del programa a mano.  Así el mensaje sigue siendo correcto aunque
  renombres el ejecutable.
\end{consejo}

% =============================================================================
\section{Convertir argumentos a números}
% =============================================================================

Los argumentos siempre llegan como cadenas de texto.  Para usarlos en
operaciones numéricas hay que convertirlos.  C++ ofrece varias funciones
para esto:

\begin{center}
\begin{tabular}{lll}
  \toprule
  \textbf{Función}   & \textbf{Tipo que devuelve} & \textbf{Encabezado} \\
  \midrule
  \cpp{atoi(s)}      & \cpp{int}                  & \cpp{<cstdlib>} \\
  \cpp{atol(s)}      & \cpp{long}                 & \cpp{<cstdlib>} \\
  \cpp{atof(s)}      & \cpp{double}               & \cpp{<cstdlib>} \\
  \cpp{stoi(s)}      & \cpp{int}                  & \cpp{<string>} \\
  \cpp{stol(s)}      & \cpp{long}                 & \cpp{<string>} \\
  \cpp{stod(s)}      & \cpp{double}               & \cpp{<string>} \\
  \bottomrule
\end{tabular}
\end{center}

Las funciones que comienzan con \texttt{a} (de \textit{ASCII to}) vienen
de~C y reciben \cpp{char*}.  Las que comienzan con \texttt{s} (de
\textit{string to}) son de C++ y aceptan \cpp{std::string}.  Como
\cpp{argv[i]} ya es \cpp{char*}, las versiones de~C funcionan
directamente sin conversión extra.

\begin{advertencia}
  \cpp{atoi} devuelve~\texttt{0} cuando la cadena no es un número
  válido, sin ningún aviso.  \cpp{stoi} lanza una excepción en ese caso.
  Si los argumentos vienen directamente del usuario y pueden ser
  inválidos, \cpp{stoi} dentro de un bloque \cpp{try}/\cpp{catch} es
  la opción más segura.
\end{advertencia}

\begin{ejemplo}[Calculadora de dos operandos]
  \textbf{calc.cpp}:

\begin{codigo}
#include <iostream>
#include <cstdlib>
using namespace std;

int main(int argc, char* argv[]){
    if (argc != 4){
        cerr << "Uso: " << argv[0]
             << " <num1> <op> <num2>" << endl;
        return 1;
    }
    double a  = atof(argv[1]);
    char   op = argv[2][0];
    double b  = atof(argv[3]);

    if      (op == '+') cout << a + b << endl;
    else if (op == '-') cout << a - b << endl;
    else if (op == '*') cout << a * b << endl;
    else if (op == '/') cout << a / b << endl;
    else {
        cerr << "Operador desconocido: " << op << endl;
        return 1;
    }
}
\end{codigo}

  Ejecuciones:
  \begin{verbatim}
$ ./calc 3.5 + 1.5
5
$ ./calc 10 / 4
2.5
  \end{verbatim}
\end{ejemplo}

% =============================================================================
\section{Actividad: el teclado numérico}
% =============================================================================

Antes de los \textit{smartphones}, todos los teléfonos celulares tenían un
teclado de doce teclas.  Para escribir un mensaje de texto había que
presionar cada tecla varias veces: el número de pulsaciones indicaba qué
letra se quería.

El detalle importante es que la \textbf{primera pulsación de una tecla
siempre introduce el dígito}, no una letra.  Las letras se obtienen a
partir de la segunda pulsación.  Con el mapeo estándar (ITU E.161, 1994):

\begin{center}
\begin{tabular}{cll}
  \toprule
  \textbf{Tecla} & \textbf{Letras} & \textbf{Pulsaciones para cada letra} \\
  \midrule
  2 & a b c & \texttt{22} \quad \texttt{222} \quad \texttt{2222} \\
  3 & d e f & \texttt{33} \quad \texttt{333} \quad \texttt{3333} \\
  4 & g h i & \texttt{44} \quad \texttt{444} \quad \texttt{4444} \\
  5 & j k l & \texttt{55} \quad \texttt{555} \quad \texttt{5555} \\
  6 & m n o & \texttt{66} \quad \texttt{666} \quad \texttt{6666} \\
  7 & p q r s & \texttt{77} \quad \texttt{777} \quad \texttt{7777} \quad \texttt{77777} \\
  8 & t u v & \texttt{88} \quad \texttt{888} \quad \texttt{8888} \\
  9 & w x y z & \texttt{99} \quad \texttt{999} \quad \texttt{9999} \quad \texttt{99999} \\
  \bottomrule
\end{tabular}
\end{center}

La regla general: para la $n$-ésima letra de una tecla se presiona ese
dígito exactamente $n+1$ veces (una por el dígito más $n$ por la posición).
Los espacios entre palabras se representan con un espacio en la salida.

\subsection*{Ejemplo verificado}

\begin{verbatim}
$ ./convertir hola luna
4446666555522 555588866622
\end{verbatim}

Comprobación de \texttt{hola}:

\begin{center}
\begin{tabular}{clccc}
  \toprule
  \textbf{Letra} & \textbf{Tecla} & \textbf{Posición} & \textbf{Pulsaciones} & \textbf{Salida} \\
  \midrule
  h & 4 (ghi) & 2ª & $2+1=3$ & \texttt{444} \\
  o & 6 (mno) & 3ª & $3+1=4$ & \texttt{6666} \\
  l & 5 (jkl) & 3ª & $3+1=4$ & \texttt{5555} \\
  a & 2 (abc) & 1ª & $1+1=2$ & \texttt{22} \\
  \bottomrule
\end{tabular}
\end{center}

\begin{center}
  \texttt{hola} $\longrightarrow$ \texttt{444} + \texttt{6666} +
  \texttt{5555} + \texttt{22} = \texttt{4446666555522}
\end{center}

Comprobación de \texttt{luna}:

\begin{center}
\begin{tabular}{clccc}
  \toprule
  \textbf{Letra} & \textbf{Tecla} & \textbf{Posición} & \textbf{Pulsaciones} & \textbf{Salida} \\
  \midrule
  l & 5 (jkl) & 3ª & $3+1=4$ & \texttt{5555} \\
  u & 8 (tuv) & 2ª & $2+1=3$ & \texttt{888} \\
  n & 6 (mno) & 2ª & $2+1=3$ & \texttt{666} \\
  a & 2 (abc) & 1ª & $1+1=2$ & \texttt{22} \\
  \bottomrule
\end{tabular}
\end{center}

\begin{center}
  \texttt{luna} $\longrightarrow$ \texttt{5555} + \texttt{888} +
  \texttt{666} + \texttt{22} = \texttt{555588866622}
\end{center}

\subsection*{Especificación del programa}

\begin{itemize}
  \item El programa se llama \texttt{convertir} y recibe las palabras
        como argumentos de consola (cada palabra es un \cpp{argv[i]}).
  \item Cada argumento contiene solo letras minúsculas del alfabeto inglés.
  \item La salida imprime las pulsaciones de cada palabra separadas por
        un espacio, seguidas de un salto de línea.
  \item Para simplificar, ningún argumento contendrá dos letras
        consecutivas en la misma tecla (por ejemplo, no habrá
        \texttt{"gh"} ni \texttt{"mn"} juntas dentro de la misma palabra).
\end{itemize}

\subsection*{Pista}

Observa que las letras del alfabeto inglés van de \cpp{'a'} (ASCII~97)
a \cpp{'z'} (ASCII~122).  La expresión \cpp{letra - 'a'} produce un
índice de~0 a~25 que puedes usar para buscar en una tabla.  Piensa en
qué datos necesitas guardar por cada letra: su tecla y su posición en
esa tecla.

% =============================================================================
\section{Ejercicios}
% =============================================================================

\begin{enumerate}

  \item Escribe un programa \textbf{eco.cpp} que imprima cada argumento
        en una línea separada, numerado desde~1.  Ejemplo:
        \begin{verbatim}
$ ./eco rojo verde azul
1: rojo
2: verde
3: azul
        \end{verbatim}
        Si no se proporcionan argumentos, el programa no imprime nada
        (sin mensaje de error).

  \item Escribe un programa \textbf{mayor.cpp} que reciba varios enteros
        como argumentos e imprima el mayor de ellos.  Si no se
        proporciona ningún argumento, muestra un mensaje de uso en
        \cpp{cerr} y termina con código~1.

  \item Escribe un programa \textbf{reversa.cpp} que imprima los
        argumentos en orden inverso (el último primero, el primero al
        final).  Ejemplo:
        \begin{verbatim}
$ ./reversa uno dos tres
tres
dos
uno
        \end{verbatim}

  \item Escribe un programa \textbf{contador.cpp} que reciba como primer
        argumento una letra y como argumentos siguientes varias palabras,
        y cuente cuántas veces aparece esa letra en el conjunto de palabras.
        Para recorrer cada argumento usa el patrón de \texttt{argumentos4.cpp}:
        itera carácter a carácter hasta encontrar \texttt{'\textbackslash{}0'}.
        Ejemplo:
        \begin{verbatim}
$ ./contador a banana manzana naranja
7
        \end{verbatim}

  \item Escribe el programa \textbf{convertir.cpp} de la actividad.
        Antes de entregarlo, incluye al inicio del archivo, como
        comentarios, los cálculos manuales de al menos dos palabras de
        prueba distintas que usaste para verificar tu solución.

\end{enumerate}

% =============================================================================
\section*{Resumen}
% =============================================================================

\begin{itemize}
  \item \cpp{main} puede recibir dos parámetros: \cpp{int argc}
        (número de cadenas, incluyendo el nombre del programa) y
        \cpp{char* argv[]} (arreglo de cadenas~C con cada argumento).
  \item Las notaciones \cpp{char* argv[]} y \cpp{char** argv} son
        equivalentes; ambas declaran un apuntador a apuntador a \cpp{char}.
  \item \cpp{argv[0]} siempre es el nombre del ejecutable; los argumentos
        del usuario están en \cpp{argv[1]} a \cpp{argv[argc-1]}.
  \item Siempre verifica el valor de \cpp{argc} antes de acceder a
        \cpp{argv[i]}.  Leer fuera de rango es comportamiento indefinido.
  \item Cada \cpp{argv[i]} es un \cpp{char*}: una cadena de C terminada
        en \texttt{'\textbackslash{}0'} que puedes recorrer carácter a carácter.
  \item Los argumentos llegan como cadenas; \cpp{atoi}, \cpp{atof} y
        \cpp{stoi} permiten convertirlos a enteros o flotantes.
  \item El mensaje de uso convencional se imprime en \cpp{cerr} y el
        programa retorna~\texttt{1} para señalar error al shell.
  \item En el teclado numérico telefónico, la $n$-ésima letra de una
        tecla requiere $n+1$ pulsaciones: una por el dígito y $n$ por
        su posición en la tecla.
  \item La filosofía Unix de ``hacer una cosa y hacerla bien'' se apoya
        en los argumentos de consola: cada programa recibe su
        configuración por \cpp{argv} y puede encadenarse con otros
        mediante tuberías.
\end{itemize}

\end{document}
